\abstract{РЕФЕРАТ}

Объем работы равен \formbytotal{lastpage}{страниц}{е}{ам}{ам}. Работа содержит \formbytotal{figurecnt}{иллюстраци}{ю}{и}{й}, \formbytotal{tablecnt}{таблиц}{у}{ы}{}, \arabic{bibcount} библиографических источников и \formbytotal{числоПлакатов}{лист}{}{а}{ов} графического материала. Количество приложений – 2. Графический материал представлен в приложении А. Фрагменты исходного кода представлены в приложении Б.

Перечень ключевых слов: городской гид, веб-приложение, Курск, интерактивная карта, REST API, Node.js, Express, PostgreSQL, Yandex Maps API, администратор, маршруты, достопримечательности, фильтрация, CRUD, фотогалерея, цифровая платформа, клиент-сервер.

Объектом разработки является веб-приложение «Цифровая платформа Курский край» — интерактивный туристический гид по городу Курску.

Целью выпускной квалификационной работы является создание сервиса для жителей и гостей города, позволяющего просматривать достопримечательности, строить маршруты, получать информацию об объектах, а также предоставляющего администратору инструменты для управления контентом.

В процессе создания сайта разработаны:
\begin{itemize}
	\item клиентская часть на HTML/CSS/JavaScript с использованием Yandex Maps API;
	\item серверная часть на Node.js (Express), предоставляющая REST API;
	\item база данных PostgreSQL, содержащая таблицы категорий, мест, маршрутов и точек маршрутов;
	\item административная панель для CRUD-операций над объектами и маршрутами;
	\item модуль построения пешеходных маршрутов с визуализацией на карте;
	\item фотогалерея с привязкой к объектам.
\end{itemize}

При разработке использовались современные веб-технологии, обеспечивающие адаптивность, производительность и удобство интерфейса.

\selectlanguage{english}
\abstract{ABSTRACT}
  
The volume of work is \formbytotal{lastpage}{page}{}{s}{s}. The work contains \formbytotal{figurecnt}{illustration}{}{s}{s}, \formbytotal{tablecnt}{table}{}{s}{s}, \arabic{bibcount} bibliographic sources and \formbytotal{числоПлакатов}{sheet}{}{s}{s} of graphic material. The number of applications is 2. The graphic material is presented in annex A. The layout of the site, including the connection of components, is presented in annex B.

List of keywords: city guide, web application, Kursk, interactive map, REST API, Node.js, Express, PostgreSQL, Yandex Maps API, administration, routes, attractions, filtering, CRUD, photo gallery, digital platform, client-server.

The object of development is the web application “Kursk Krai Digital Platform” — an interactive tourist guide to the city of Kursk.

The purpose of the final qualifying work is to create a convenient service for residents and guests of the city, allowing them to view attractions, build routes, receive information about objects, as well as provide the administrator with tools for content management.

During the creation of the site, the following were developed:
\begin{itemize}
	\item client-side in HTML/CSS/JavaScript using Yandex Maps API;
	\item server-side in Node.js (Express) providing REST API;
	\item PostgreSQL database containing tables for categories, places, routes, and route points;
	\item admin panel for CRUD operations on objects and routes;
	\item pedestrian route construction module with map visualization;
	\item photo gallery linked to objects.
\end{itemize}

Modern web technologies were used to ensure responsiveness, performance, and user-friendly interface.
\selectlanguage{russian}
